\FloatBarrier
\section{Esercitazioni Guidate sul Classificatore \texorpdfstring{$K$-NN}{K-NN}}
\label{sec:dm-knn-exercises}

\FloatBarrier
\subsection{Esercizio 1: Classificazione su Griglia Bidimensionale (\texorpdfstring{$k = 3$}{k = 3})}
\label{subsec:dm-knn-ex1}

Si consideri un insieme di addestramento su griglia cartesiana composto da $12$ elementi numerati da $1$ a $12$, etichettati in due classi: \texttt{CIRCLE} (cerchi) e \texttt{DIAMOND} (rombi). Si richiede di classificare tre istanze di test $A$, $B$ e $C$ (non appartenenti al training set) mediante $K$-NN con $k = 3$ e distanza Euclidea:
\begin{enumerate}
    \item \textbf{Punto di test $A$:} i $3$ vicini pi\`u prossimi risultano $k\text{NN}(A) = \{3, 5, 12\}$, tutti di classe \texttt{CIRCLE}. Con $3$ voti a $0$, $A$ viene classificato come \texttt{CIRCLE}.
    \item \textbf{Punto di test $B$:} i due vicini strettamente pi\`u prossimi sono $3$ e $5$ (\texttt{CIRCLE}), mentre i punti $7$ (\texttt{CIRCLE}) e $10$ (\texttt{DIAMOND}) si trovano esattamente alla \textbf{stessa distanza Euclidea} da $B$ (situazione di pareggio o \textit{tie} sulla terza posizione). Qualunque sia il criterio di spareggio adottato tra $7$ e $10$, l'insieme dei $3$ vicini conterr\`a $\{3, 5, 7\}$ ($3$ \texttt{CIRCLE}) oppure $\{3, 5, 10\}$ ($2$ \texttt{CIRCLE} e $1$ \texttt{DIAMOND}). In entrambi i casi $B$ viene classificato come \texttt{CIRCLE}.
    \item \textbf{Punto di test $C$:} la circonferenza centrata in $C$ racchiude i $3$ vicini $k\text{NN}(C) = \{4, 6, 7\}$, dove $6$ e $7$ sono \texttt{CIRCLE} mentre $4$ \`e \texttt{DIAMOND} (il punto $1$ giace appena all'esterno). Con $2$ voti contro $1$, $C$ viene classificato come \texttt{CIRCLE}.
\end{enumerate}

\FloatBarrier
\subsection{Esercizio 2: Svolgimento Analitico su Griglia Cartesiana (\texorpdfstring{$k = 3$}{k = 3})}
\label{subsec:dm-knn-ex2}

Su una seconda griglia cartesiana bidimensionale (\autoref{fig:dm-knn-grid-ex2}), fissando l'origine $(0,0)$ nell'intersezione in basso a sinistra, sono collocati $12$ punti di training su coordinate intere esatte:
\begin{itemize}
    \item Classe \texttt{CIRCLE} (blu): $1=(14,8)$, $3=(3,1)$, $5=(4,4)$, $6=(7,6)$, $7=(6,5)$, $9=(2,8)$, $12=(6,7)$.
    \item Classe \texttt{DIAMOND} (rosso): $2=(11,5)$, $4=(9,8)$, $8=(12,8)$, $10=(2,6)$, $11=(4,9)$.
\end{itemize}
Si classificano i tre punti di test $A=(8,3)$, $B=(10,3)$ e $C=(12,3)$ utilizzando $k = 3$ e la distanza Euclidea $d(\mathbf{p}, \mathbf{q}) = \sqrt{(p_x - q_x)^2 + (p_y - q_y)^2}$:

\begin{enumerate}
    \item \textbf{Classificazione di $A = (8,3)$:}
    \begin{align*}
        d(A, 7) &= \sqrt{(8-6)^2 + (3-5)^2} = \sqrt{8} \approx 2.83 \; (\text{CIRCLE}) \\
        d(A, 6) &= \sqrt{(8-7)^2 + (3-6)^2} = \sqrt{10} \approx 3.16 \; (\text{CIRCLE}) \\
        d(A, 2) &= \sqrt{(8-11)^2 + (3-5)^2} = \sqrt{13} \approx 3.61 \; (\text{DIAMOND})
    \end{align*}
    Tutti gli altri punti hanno distanza $\ge \sqrt{17} \approx 4.12$. Quindi $k\text{NN}(A) = \{7, 6, 2\}$ ($2$ \texttt{CIRCLE}, $1$ \texttt{DIAMOND}) $\implies$ \textbf{$A$ \`e classificato come \texttt{CIRCLE}}.

    \item \textbf{Classificazione di $B = (10,3)$:}
    \begin{align*}
        d(B, 2) &= \sqrt{(10-11)^2 + (3-5)^2} = \sqrt{5} \approx 2.24 \; (\text{DIAMOND}) \\
        d(B, 6) &= \sqrt{(10-7)^2 + (3-6)^2} = \sqrt{18} \approx 4.24 \; (\text{CIRCLE}) \\
        d(B, 7) &= \sqrt{(10-6)^2 + (3-5)^2} = \sqrt{20} \approx 4.47 \; (\text{CIRCLE})
    \end{align*}
    I successivi pi\`u vicini sono $4$ ($\sqrt{26} \approx 5.10$) e $8$ ($\sqrt{29} \approx 5.39$). Quindi $k\text{NN}(B) = \{2, 6, 7\}$ ($2$ \texttt{CIRCLE}, $1$ \texttt{DIAMOND}) $\implies$ \textbf{$B$ \`e classificato come \texttt{CIRCLE}}.

    \item \textbf{Classificazione di $C = (12,3)$:}
    \begin{align*}
        d(C, 2) &= \sqrt{(12-11)^2 + (3-5)^2} = \sqrt{5} \approx 2.24 \; (\text{DIAMOND}) \\
        d(C, 8) &= \sqrt{(12-12)^2 + (3-8)^2} = \sqrt{25} = 5.00 \; (\text{DIAMOND}) \\
        d(C, 1) &= \sqrt{(12-14)^2 + (3-8)^2} = \sqrt{29} \approx 5.39 \; (\text{CIRCLE})
    \end{align*}
    I successivi pi\`u vicini sono $4$ e $6$ (entrambi a distanza $\sqrt{34} \approx 5.83$). Quindi $k\text{NN}(C) = \{2, 8, 1\}$ ($2$ \texttt{DIAMOND}, $1$ \texttt{CIRCLE}) $\implies$ \textbf{$C$ \`e classificato come \texttt{DIAMOND}}.
\end{enumerate}

\begin{figure}[H]
\centering
\resizebox{0.92\columnwidth}{!}{%
\begin{tikzpicture}[scale=0.55]
    % Griglia
    \draw[step=1cm, gray!30, very thin] (0,0) grid (15,10);
    \draw[->, gray!70] (0,0) -- (15.4,0) node[right] {$x$};
    \draw[->, gray!70] (0,0) -- (0,10.4) node[above] {$y$};

    % Punti CIRCLE (blu)
    \foreach \x/\y/\lbl in {14/8/1, 3/1/3, 4/4/5, 7/6/6, 6/5/7, 2/8/9, 6/7/12} {
        \fill[navy] (\x,\y) circle (5pt);
        \node[above right, font=\scriptsize\bfseries, color=navy] at (\x,\y) {\lbl};
    }

    % Punti DIAMOND (rosso)
    \foreach \x/\y/\lbl in {11/5/2, 9/8/4, 12/8/8, 2/6/10, 4/9/11} {
        \node[diamond, fill=crimson, inner sep=2.5pt] at (\x,\y) {};
        \node[above right, font=\scriptsize\bfseries, color=crimson] at (\x,\y) {\lbl};
    }

    % Punti di test (verde)
    \foreach \x/\y/\lbl in {8/3/A, 10/3/B, 12/3/C} {
        \node[font=\bfseries\small, color=darkgreen] at (\x,\y) {$\times$};
        \node[below, font=\scriptsize\bfseries, color=darkgreen] at (\x,\y-0.1) {\lbl};
    }

    % Intorno di A (raggio sqrt(13) ~ 3.61) e C (raggio sqrt(29) ~ 5.39)
    \draw[dashed, navy!80, line width=0.8pt] (8,3) circle (3.65);
    \draw[dotted, crimson!80, line width=0.9pt] (12,3) circle (5.42);
\end{tikzpicture}%
}
\caption{Rappresentazione cartesiana dell'Esercizio 2: cerchi pieni (`\texttt{CIRCLE}'), rombi (`\texttt{DIAMOND}') e punti di test $A, B, C$ con gli intorni di $A$ (tratteggiato) e $C$ (punteggiato).}
\label{fig:dm-knn-grid-ex2}
\end{figure}

\FloatBarrier
\subsection{Esercizio 3: Case-Based Reasoning Medico con Similarit\`a Asimmetrica Pesata}
\label{subsec:dm-knn-ex3}

Un sistema di \textit{Case-Based Reasoning} per la diagnosi medica memorizza $7$ pazienti ($c_1, \dots, c_7$) descritti da quattro attributi sintomatologici (\texttt{Fever} $F$, \texttt{Vomiting} $V$, \texttt{Diarrhea} $D$, \texttt{Shivering} $Sh$) e da una diagnosi clinica (\autoref{tab:dm-cbr-cases}).

\begin{table}[H]
\centering
\caption{Base di casi medici ($c_1, \dots, c_7$) per il sistema di Case-Based Reasoning.}
\label{tab:dm-cbr-cases}
\small
\resizebox{\columnwidth}{!}{%
\begin{tabular}{cccccc}
\toprule
\textbf{Case} & \textbf{Fever ($F$)} & \textbf{Vomiting ($V$)} & \textbf{Diarrhea ($D$)} & \textbf{Shivering ($Sh$)} & \textbf{Classification} \\
\midrule
$c_1$ & no      & no  & no  & no  & healthy (H) \\
$c_2$ & average & no  & no  & no  & influenza (I) \\
$c_3$ & high    & no  & no  & yes & influenza (I) \\
$c_4$ & high    & yes & yes & no  & salmonella poisoning (S) \\
$c_5$ & average & no  & yes & no  & salmonella poisoning (S) \\
$c_6$ & no      & yes & yes & no  & bowel inflammation (B) \\
$c_7$ & average & yes & yes & no  & bowel inflammation (B) \\
\bottomrule
\end{tabular}%
}
\end{table}

Le matrici di similarit\`a locale fornite dall'esperto sono \textbf{asimmetriche} (le righe rappresentano il valore nella query $q$, le colonne il valore nel caso memorizzato $c$) e sono riportate nella \autoref{tab:dm-cbr-sim}.

\begin{table}[H]
\centering
\caption{Matrici di similarit\`a locale asimmetrica per \texttt{Fever} ($\text{sim}_F$) e per \texttt{Vomiting}, \texttt{Diarrhea}, \texttt{Shivering} ($\text{sim}_V = \text{sim}_D = \text{sim}_{Sh}$).}
\label{tab:dm-cbr-sim}
\small
\begin{tabular}{lccc@{\hspace{1.4cm}}lcc}
\toprule
$\text{sim}_F(q, c)$ & \textbf{no} & \textbf{avg} & \textbf{high} & $\text{sim}_{V,D,Sh}(q, c)$ & \textbf{yes} & \textbf{no} \\
\midrule
\textbf{no}   & $1.0$ & $0.7$ & $0.2$ & \textbf{yes} & $1.0$ & $0.0$ \\
\textbf{avg}  & $0.5$ & $1.0$ & $0.8$ & \textbf{no}  & $0.2$ & $1.0$ \\
\textbf{high} & $0.0$ & $0.3$ & $1.0$ & & & \\
\bottomrule
\end{tabular}
\end{table}

I pesi degli attributi sono $w_F = 0.3$, $w_V = 0.2$, $w_D = 0.2$ e $w_{Sh} = 0.3$, e la similarit\`a globale \`e definita come:
\begin{equation}
\begin{split}
    \text{Sim}(q, c) &= w_F \cdot \text{sim}_F(q_F, c_F) + w_V \cdot \text{sim}_V(q_V, c_V) \\
    &\quad + w_D \cdot \text{sim}_D(q_D, c_D) + w_{Sh} \cdot \text{sim}_{Sh}(q_{Sh}, c_{Sh})
\end{split}
\end{equation}

Si vuole classificare la nuova istanza di query $q = (\text{high}, \text{no}, \text{no}, \text{no})$ per $K = 1, 2, 3$ selezionando i $K$ casi con similarit\`a globale maggiore:
\begin{align*}
    \text{Sim}(q, c_1) &= 0.3(0.0) + 0.2(1.0) + 0.2(1.0) + 0.3(1.0) = 0.70 \\
    \text{Sim}(q, c_2) &= 0.3(0.3) + 0.2(1.0) + 0.2(1.0) + 0.3(1.0) = 0.79 \\
    \text{Sim}(q, c_3) &= 0.3(1.0) + 0.2(1.0) + 0.2(1.0) + 0.3(0.2) = 0.76 \\
    \text{Sim}(q, c_4) &= 0.3(1.0) + 0.2(0.2) + 0.2(0.2) + 0.3(1.0) = 0.68 \\
    \text{Sim}(q, c_5) &= 0.3(0.3) + 0.2(1.0) + 0.2(0.2) + 0.3(1.0) = 0.63 \\
    \text{Sim}(q, c_6) &= 0.3(0.0) + 0.2(0.2) + 0.2(0.2) + 0.3(1.0) = 0.38 \\
    \text{Sim}(q, c_7) &= 0.3(0.3) + 0.2(0.2) + 0.2(0.2) + 0.3(1.0) = 0.47
\end{align*}

Ordinando i casi per similarit\`a decrescente si ottiene la graduatoria $c_2 (0.79) > c_3 (0.76) > c_1 (0.70) > c_4 (0.68) > c_5 (0.63) > c_7 (0.47) > c_6 (0.38)$. Pertanto:
\begin{itemize}
    \item \textbf{Per $K = 1$:} si seleziona $c_2$ ($\text{Sim} = 0.79$) $\implies$ \textbf{Classe predetta: \texttt{Influenza}}.
    \item \textbf{Per $K = 2$:} si selezionano $c_2$ (\texttt{Influenza}) e $c_3$ (\texttt{Influenza}) $\implies$ \textbf{Classe predetta: \texttt{Influenza}} (unanimit\`a $2/2$).
    \item \textbf{Per $K = 3$:} si selezionano $c_2$ (\texttt{Influenza}), $c_3$ (\texttt{Influenza}) e $c_1$ (\texttt{healthy}) $\implies$ con $2$ voti a $1$, la \textbf{Classe predetta \`e \texttt{Influenza}}.
\end{itemize}

% ====================================================================
\FloatBarrier
